\documentclass[11pt]{article}
\usepackage[margin=1in]{geometry}
\usepackage{amsmath,amssymb}
\usepackage{booktabs}
\usepackage{siunitx}
\usepackage{xcolor}
\usepackage[colorlinks=true,linkcolor=blue,urlcolor=blue]{hyperref}

\newcommand{\htrue}{h_{\mathrm{true}}}
\newcommand{\hhat}{\hat h}
\newcommand{\dz}{\sigma_z}
\newcommand{\dw}{\sigma_w}
\newcommand{\rhohat}{\hat\rho}
\newcommand{\Cov}{\mathrm{Cov}}

\title{Why the diagonal--Gaussian estimator of the dark--matter scale height\\
sits about $1\sigma$ away from the truth}
\author{Phase--space inference notes}
\date{\today}

\begin{document}
\maketitle

\begin{abstract}
When we recover the dark--matter scale height $h_{\rm dm}$ by back--integrating an
observed phase--space snapshot and fitting a diagonal Gaussian (equivalently, a
linear ``$a u + b$'' bijection) to the initial cloud, the profile--likelihood
minimum lands a few parsec away from the true value---for our fiducial
$N=\num{20000}$, seed--7 data set it sits at $h\simeq\SI{258}{pc}$ versus a truth
of \SI{250}{pc}. This note shows that the offset is \emph{not} a systematic
bias, an integrator artifact, or a modelling error: it is the ordinary
$\sim\!1\sigma$ statistical scatter of the estimator. We derive the error bar
$\sigma_h=\sigma_\rho/|d\rho/dh|\approx(1/\sqrt N)/|d\rho/dh|$ by propagating the
sampling noise of a correlation coefficient through the estimating equation, and
verify it reproduces the measured seed--to--seed scatter of \SI{8}{pc}.
\end{abstract}

\section{The observation}
The data are $N$ stars drawn from a centred, separable initial distribution
function $f_0(z_0,w_0)=\mathcal N(0,\dz^{\,0})\,\mathcal N(0,\dw^{\,0})$, evolved
for $t_{\rm obs}=\SI{400}{Myr}$ in the true potential (scale height $\htrue=\SI{250}{pc}$),
with no satellite. To infer $h$ we back--integrate the snapshot to $t=0$ under a
trial potential of scale height $h$, obtaining a cloud $(z_0,w_0)=T_h(z,w)$, and
fit a \emph{diagonal} Gaussian to it. The negative log--likelihood of that model
depends only on the marginal widths,
\begin{equation}
  \mathrm{NLL}(h) \;=\; N\Big[\ln(2\pi) + \ln\dz(h) + \ln\dw(h) + 1\Big],
  \label{eq:nll}
\end{equation}
so the estimator $\hhat=\arg\min_h \mathrm{NLL}(h)$ minimises the product
$\dz(h)\,\dw(h)$. Empirically the minimum sits a few pc above $\htrue$, and it
does so \emph{reproducibly} across our tests---simply because those tests reuse
the same random seed.

\section{What the estimator really measures}
Write the $2\times2$ covariance of the back--integrated cloud as
\begin{equation}
  \Cov(h) =
  \begin{pmatrix} \dz^2 & \rho\,\dz\dw \\ \rho\,\dz\dw & \dw^2 \end{pmatrix},
  \qquad
  \det\Cov(h) = \dz^2\dw^2\,(1-\rho^2),
\end{equation}
where $\rho(h)$ is the $z$--$w$ correlation coefficient. Hence
\begin{equation}
  \ln\big(\dz\dw\big) = \tfrac12\ln\det\Cov(h) - \tfrac12\ln\!\big(1-\rho(h)^2\big).
  \label{eq:decomp}
\end{equation}
The back--integration $T_h$ is a \emph{symplectic} map, so it conserves the
phase--space measure. For these near--harmonic clouds the second--moment area
$\det\Cov(h)$ is conserved to excellent accuracy as well (numerically it is flat
to $<0.1\%$ across the scanned range, $\det\Cov\simeq\num{2.645e7}$). With the
first term in \eqref{eq:decomp} constant, minimising $\ln(\dz\dw)$ is
\emph{equivalent to} maximising $1-\rho^2$, i.e.\ to the condition
\begin{equation}
  \boxed{\;\rhohat(\hhat)=0\;}
  \label{eq:estimating}
\end{equation}
The diagonal model, which ignores the correlation, therefore places its optimum
exactly at the scale height for which the recovered cloud is \emph{uncorrelated}.

\section{Why $\rho=0$ is not exactly at the truth}
At the true scale height the map $T_{\htrue}$ inverts the forward evolution and
returns the original cloud, whose \emph{population} correlation is zero. But the
\emph{sample} correlation of a finite draw is not exactly zero: for $N$ points
with true correlation $\approx0$ it fluctuates about $0$ with a spread
$\sigma_\rho$. Thus
\begin{equation}
  \rhohat(\htrue) \;=\; \underbrace{\rho_{\rm true}(\htrue)}_{=\,0} + \;\delta\rho,
  \qquad \mathbb E[\delta\rho]=0,\quad \mathrm{SD}[\delta\rho]=\sigma_\rho .
\end{equation}
Because $\rhohat(\htrue)\neq0$ in any single realisation, the zero of
$\rhohat(h)$---and hence $\hhat$---is displaced from $\htrue$.

\section{The error bar: propagating $\sigma_\rho$ through the estimating equation}
This is the delta method (equivalently, the implicit--function theorem) applied
to \eqref{eq:estimating}. Linearise $\rhohat(h)$ about the truth,
\begin{equation}
  \rhohat(h) \;\approx\; \rhohat(\htrue) + \frac{d\rho}{dh}\,(h-\htrue),
\end{equation}
and impose $\rhohat(\hhat)=0$:
\begin{equation}
  \frac{d\rho}{dh}\,(\hhat-\htrue) + \delta\rho \approx 0
  \qquad\Longrightarrow\qquad
  \hhat-\htrue \;\approx\; -\,\frac{\delta\rho}{\,d\rho/dh\,}.
  \label{eq:offset}
\end{equation}
Taking the standard deviation of \eqref{eq:offset} gives the central result,
\begin{equation}
  \boxed{\;\sigma_h \;=\; \frac{\sigma_\rho}{\left|d\rho/dh\right|}\;}
  \label{eq:sigmah}
\end{equation}
This is the generic standard error of a root of an estimating equation
$g(\hat\theta)=0$, namely $\mathrm{SE}(\hat\theta)=\mathrm{SD}[g]_{\rm truth}/|g'|$;
it is the same ``output error $=$ input error $/$ sensitivity'' one uses for any
indirect measurement.

\paragraph{The $1/\sqrt N$ factor.} The remaining ingredient is the sampling
scatter of a correlation coefficient. Fisher's classical result for the sample
correlation $\hat\rho$ of a bivariate normal sample is
\begin{equation}
  \mathrm{Var}(\hat\rho) \;\approx\; \frac{(1-\rho^2)^2}{N}
  \;\xrightarrow[\;\rho\to0\;]{}\; \frac1N,
  \qquad\text{so}\qquad \sigma_\rho \approx \frac{1}{\sqrt N}.
\end{equation}
Substituting into \eqref{eq:sigmah},
\begin{equation}
  \sigma_h \;\approx\; \frac{1/\sqrt N}{\left|d\rho/dh\right|}.
  \label{eq:sigmah_N}
\end{equation}
$N$ enters only here: more stars pin $\rho$ tighter, but the small slope
$d\rho/dh$---the weak, anharmonic grip of $h$ on the data---sets a large lever
arm that inflates the resulting uncertainty in $h$.

\section{Numerical verification}
For the fiducial data ($N=\num{20000}$, forward integration with
\num{12000} steps) the decomposition \eqref{eq:decomp} is confirmed directly:
$\det\Cov$ is flat, $\rho(h)$ marches monotonically through zero, and the
$\ln(\dz\dw)$ minimum coincides with the $\rho=0$ crossing
(Table~\ref{tab:decomp}).

\begin{table}[h]
\centering
\begin{tabular}{rrrr}
\toprule
$h$ [pc] & $\dz$ [pc] & $\rho$ & $\det\Cov$ \\
\midrule
230 & 258.1 & $+0.033$ & \num{2.646e7} \\
250 & 257.4 & $+0.010$ & \num{2.645e7} \\
259 & 257.2 & $\approx 0$ & \num{2.645e7} \\
270 & 257.2 & $-0.011$ & \num{2.645e7} \\
\bottomrule
\end{tabular}
\caption{The correlation $\rho(h)$ passes through $0$ near $h\simeq\SI{259}{pc}$,
while $\det\Cov$ is essentially constant. The diagonal--Gaussian minimum tracks
the $\rho=0$ point, not the truth.}
\label{tab:decomp}
\end{table}

Repeating the fit over independent random seeds shows the offset is zero--mean and
of the size predicted by \eqref{eq:sigmah} (Table~\ref{tab:seeds}). The measured
slope is $|d\rho/dh|\approx\SI{0.0011}{pc^{-1}}$, and the sample scatter of the
initial--cloud correlation across seeds is $\sigma_\rho=0.0088$ (versus the ideal
$1/\sqrt{\num{20000}}=0.0071$; the excess comes from the evolve/back--integrate
round trip mildly amplifying the noise). Equation~\eqref{eq:sigmah} then predicts
\begin{equation}
  \sigma_h \;=\; \frac{0.0088}{0.0011} \;\approx\; \SI{8.0}{pc},
\end{equation}
in close agreement with the directly measured seed--to--seed scatter of
\SI{8.4}{pc}. The fiducial seed--7 result, $\hhat=\SI{259}{pc}$, is thus
$+1.1\sigma$---an entirely ordinary draw.

\begin{table}[h]
\centering
\begin{tabular}{lrrr}
\toprule
& $\rho_{\rm IC}$ & $\hhat-\htrue$ [pc] & \\
\midrule
seed 7 (fiducial) & $+0.0100$ & $+9.4$ & \\
seed 1 & $+0.0036$ & $+3.7$ & \\
seed 3 & $-0.0012$ & $-0.9$ & \\
seed 2 & $-0.0130$ & $-11.9$ & \\
seed 6 & $-0.0130$ & $-12.6$ & \\
\midrule
mean ($8$ seeds) & $\approx 0$ & $-3.5$ & (consistent with $0$) \\
scatter ($8$ seeds) & $0.0088$ & $8.4$ & $\;\approx \sigma_h$ \\
$N=\num{80000}$ scatter & $0.0038$ & $2.4$ & $\;\propto 1/\sqrt N$ \\
\bottomrule
\end{tabular}
\caption{The minimum offset tracks the finite--sample initial--cloud correlation
$\rho_{\rm IC}$, averages to zero, and shrinks as $1/\sqrt N$ (scatter
$\SI{8.4}{pc}\to\SI{2.4}{pc}$ from $N=\num{20000}$ to $\num{80000}$).}
\label{tab:seeds}
\end{table}

\section{Conclusion}
The few--parsec displacement of the profile minimum is the $1\sigma$ statistical
error of the estimator, not a bias:
\begin{enumerate}\itemsep2pt
  \item the diagonal/affine model reads $h$ only through the condition
    $\rhohat(h)=0$;
  \item in a finite sample $\rhohat(\htrue)\neq0$ by $\mathcal O(1/\sqrt N)$;
  \item the shallow slope $d\rho/dh$ converts that tiny correlation error into a
    displacement $\sigma_h=\sigma_\rho/|d\rho/dh|\approx(1/\sqrt N)/|d\rho/dh|$.
\end{enumerate}
For our data this is $\approx\SI{8}{pc}$ at $N=\num{20000}$, so the fiducial
$\hhat=\SI{258}{pc}$ should be read as $h=250\pm8$~pc. The uncertainty is large
not because the data are few but because each star carries little information
about $h$ (the anharmonic insensitivity encoded in the small $d\rho/dh$); a
richer estimator that uses more than the correlation could tighten it.
\end{document}
